\documentclass[10pt, a4paper, twoside]{article}
\usepackage{ntv-itmo}
\graphicspath{{../figures/ru/}}
\newcommand{\strategy}[1]{\texttt{#1}}

\ntvshortauthors{С.А. Постнов, А.В. Платонов, Л.К. Саськов}
\ntvshorttitle{Семантическое шардирование поискового индекса\dots}
\ntvissue{2026, том 26, № X}{2026, vol. 26, no X}

\begin{document}
\thispagestyle{ntvfirst}

\ntvrubric{Компьютерные системы и информационные технологии}{COMPUTER SCIENCE}

\ntvudc{004.912}

\ntvtitle{Метод построения распределенного поискового индекса\\на основе семантики запросов}
% TODO(Степан): окончательное название — лаконично, без сокращений.

\ntvauthors{Степан Андреевич Постнов$^{1}$, Алексей Владимирович Платонов$^{2}$, Лев Кириллович Саськов$^{3}$}
\ntvaffiliation{$^{1,2,3}$~Университет ИТМО, Санкт-Петербург, 197101, Российская Федерация}
\ntvcontact{$^{1}$~s@pstnv.ru, \tbd{ORCID}}
\ntvcontact{$^{2}$~avplatonov@itmo.ru, https://orcid.org/0000-0002-8485-1296}
\ntvcontact{$^{3}$~\tbd{e-mail, ORCID}}

\begin{ntvabstract}{Аннотация}
\abstractpart{Введение}Распределённый поисковый индекс разрезается по документам либо по словам. При разбиении по словам инвертированный список терма хранится в одном шарде целиком, документ --- в шардах всех своих термов, поэтому запрос из нескольких слов обслуживается несколькими шардами, и при случайном назначении слов их число приближается к длине запроса. Работа рассматривает размещение слов по шардам, при котором слова одного запроса попадают в один шард.
\abstractpart{Метод}Предложен метод шардирования по графу совстречаемости слов журнала запросов. Кластеризация графа задаёт разбиение словаря, вне-графовые слова присоединяются к кластерам голосованием документов, бюджет объёма ограничивает долю корпуса в одном шарде, маршрутизация выбирает опрашиваемые шарды жадным взвешенным покрытием. Реализация охватывает четыре алгоритма кластеризации и десять стратегий, включая случайные разбиения с равным числом шардов в качестве базовой линии.
\abstractpart{Основные результаты}На корпусе MS~MARCO из 241\,185 документов и 122\,063 запросов показано, что разбиения без бюджета объёма вырождаются: крупнейший шард накапливает 83,1\,\% корпуса, и ни изменение параметра кластеризации, ни 18~вариаций настройки графа не приводят его к 15-процентному бюджету. Сбалансированная стратегия при равном числе шардов превосходит случайную по всем измеренным метрикам: overlap@10 при опросе одного шарда составляет 0,780 против 0,697 (парная разность +0,083; 95\,\%-й доверительный интервал [0,077;\,0,089]), на смещённом распределении запросов прирост равен +0,109~[0,104;\,0,114]. Сформулировано правило выбора конфигурации под четыре профиля нагрузки: минимальный fanout, минимальное дублирование, равномерное распределение объёма шардов, фиксированный объём опроса.
\abstractpart{Обсуждение}Правило выбора и соответствие между ограничением эксплуатации и параметром метода не зависят от корпуса; численные значения посчитаны для одного корпуса и одного журнала запросов. Метод предполагает глобальную статистику ранжирования BM25 и дублирование порядка двадцати копий документа. Направления дальнейшей работы --- перебор бюджета объёма, реплики сбалансированной стратегии, измерение задержек вместо прокси-стоимости.
\end{ntvabstract}
\ntvkeywords{Ключевые слова}{инвертированный индекс, шардирование по словам, маршрутизация запросов, граф совстречаемости, кластеризация графов, балансировка шардов}
% TODO(Степан): 5–10 ключевых слов, главное первым; каждое должно встречаться в тексте.
\ntvacknowledgments{Благодарности}{\tbd{финансирование (организация, номер гранта) или удалить блок}}
\ntvcopyright{Постнов С.А., Платонов А.В., Саськов Л.К., 2026}

\vspace{14pt}
\begin{otherlanguage}{english}
\ntvtitle{Semantic term sharding of a search index under a shard-volume budget}
\ntvauthors{Stepan A. Postnov$^{1}$, Alexey V. Platonov$^{2}$, Lev K. Saskov$^{3}$}
\ntvaffiliation{$^{1,2,3}$~ITMO University, Saint Petersburg, 197101, Russian Federation}
\ntvcontact{$^{1}$~stepaha78@gmail.com, \tbd{ORCID}}
\ntvcontact{$^{2}$~avplatonov@itmo.ru, https://orcid.org/0000-0002-8485-1296}
\ntvcontact{$^{3}$~\tbd{e-mail, ORCID}}
\begin{ntvabstract}{Abstract}
\tbd{Translation of the Russian abstract.}
\end{ntvabstract}
\ntvkeywords{Keywords}{inverted index, term sharding, query routing, co-occurrence graph, graph clustering, shard balancing}
\ntvacknowledgments{Acknowledgements}{\tbd{funding or remove}}
\end{otherlanguage}

% ============================================================================
\begin{ntvbody}

\section{Введение}

Распределённый поисковый индекс строится двумя способами. При разбиении по документам каждый шард содержит подвыборку корпуса и весь словарь, а запрос направляется в несколько шардов, выбранных оценкой распределения релевантных документов (селективный поиск)~\cite{kulkarni2015,dai2016}. При разбиении по словам словарь разрезается между шардами: инвертированный список терма хранится целиком в одном шарде, документ --- в шардах всех своих термов. На таком хранении строятся конвейерные системы распределённого вычисления запроса~\cite{moffat2007} и крупные индексные системы с терм-распределённым хранением~\cite{risvik2013}. Разрезание словаря исключает дублирование постингов и облегчает балансировку длин списков~\cite{moffat2006}, но запрос из $m$ слов затрагивает шарды всех своих термов, и маршрутизация становится основной статьёй стоимости обращения к индексу.

Назначение слов по хешу не требует настройки, но не использует структуру нагрузки: шарды запроса почти не пересекаются, и средний запрос рассматриваемого корпуса --- 3,2 слова --- опрашивает 3,2 шарда из 132. Если размещать вместе слова, встречающиеся в одних запросах, запрос целиком обслуживается одним шардом. Источником информации о совместной встречаемости служит журнал запросов, который поисковая система накапливает в процессе работы; построение графа совстречаемости не требует разметки и дополнительного обхода корпуса.

Кластеризация графа совстречаемости --- необходимое, но не достаточное условие практичного разбиения. Документ хранится в шардах всех своих слов, поэтому кластер частых слов собирает большую часть корпуса: в проведённых измерениях крупнейший шард достигал 83,1\,\% корпуса. Разбиение с таким шардом непригодно для кластера равноправных узлов: хранение и трафик концентрируются на одном узле. В работе равномерность обеспечивается явным бюджетом объёма шарда.

Цель работы --- метод построения разбиения словаря по графу совстречаемости слов запросов с бюджетом объёма шарда. Вклад работы состоит в следующем: предложен метод шардирования с четырьмя заменяемыми алгоритмами кластеризации и его открытая реализация; разработана воспроизводимая методика оценки на реальном журнале запросов, включая сверку расчётной выдачи шардов с физическими индексами; сформулировано правило выбора конфигурации под четыре профиля нагрузки.

Далее в разделе~2 дан обзор родственных работ, в разделе~3 описан метод, в разделе~4 --- методика оценки, в разделе~5 приведены результаты, в разделе~6 обсуждаются ограничения.

\section{Связанные работы}

Селективный поиск по разбиению документов исследован в работах~\cite{kulkarni2015,dai2016}: для каждого шарда оценивается вероятность содержать релевантные документы, и запрос направляется в ограниченное подмножество шардов; разбиение может учитывать и запросную нагрузку~\cite{dai2016}. Терм-распределённое хранение рассматривалось в системе Maguro~\cite{risvik2013} и в конвейерной архитектуре вычисления запроса~\cite{moffat2007}, где запрос декомпозируется по термам и пересекается по постинговым спискам. Балансировка терм-распределённого поиска~\cite{moffat2006} перераспределяет списки между процессорами по измеренной нагрузке. В перечисленных работах критериями служат длины списков и частоты термов; объём хранимых на узле документов как ограничение, насколько известно авторам, не рассматривался. В настоящей работе именно этот объём --- первичный параметр: бюджет доли корпуса в шарде определяет число шардов и равномерность узлов.

Вес ребра графа совстречаемости --- мера ассоциации пары слов. Использован нормализованный pointwise mutual information (NPMI)~\cite{bouma2009}, устраняющий смещение PMI в пользу редких пар; патологии немаркированной PMI разобраны в~\cite{williams2022}. Альтернативные веса --- статистика $G^2$~\cite{dunning1993} и $\chi^2$; чувствительность к выбору веса и порогов графа проверена в разделе~5.

Кластеризация графа выполняется четырьмя алгоритмами, представляющими разные принципы: максимизация модулярности~\cite{newman2004} (реализация Leiden~\cite{traag2019} с гарантией связности сообществ), качество constant Potts model~\cite{reichardt2006} и его связь с пределом разрешения~\cite{traag2011}, минимизация ожидаемой длины кода случайного блуждания (Infomap)~\cite{rosvall2008} и многоуровневая минимизация сечения графа (METIS)~\cite{karypis1998}. Выбор алгоритма в методе --- заменяемый компонент; устойчивость выводов к этому выбору проверяется измерениями.

Маршрутизация запроса к шардам формулируется как взвешенное покрытие множества термов; жадный алгоритм покрытия восходит к~\cite{chvatal1979}. Ранжирование выполнено по BM25~\cite{robertson2009}. Корпус и запросы взяты из MS~MARCO~\cite{bajaj2016}; для проверки устойчивости к смещению распределения запросов --- из Natural Questions~\cite{kwiatkowski2019}.

\section{Метод}

\subsection{Постановка задачи}

Пусть $T$ --- словарь корпуса документов $\mathcal{D}$, $S$ --- множество шардов. Разбиение $\sigma$ назначает каждому терму непустое множество шардов:
\begin{equation}
	\sigma\colon T \to 2^{S}\setminus\{\varnothing\}, \quad D(d)=\bigcup_{t\in d}\sigma(t).
	\label{eq:assignment}
\end{equation}
Документ $d$ хранится в каждом шарде каждого своего терма; дублирование $\mathrm{dup}$ --- среднее число шардов документа. Запрос $q$ корректно обслуживается любым набором $S_q\subseteq S$, содержащим по шардy каждого терма запроса: объединение выдач шардов тогда воспроизводит выдачу полного индекса~\eqref{eq:invariant} при любом ранжировании, монотонном по документу.
\begin{equation}
	\mathrm{Top}_k\Bigl(\bigcup_{s\in S_q} R_s(q)\Bigr)=\mathrm{Top}_k\bigl(R(q)\bigr).
	\label{eq:invariant}
\end{equation}

Стоимость запроса складывается из числа опрашиваемых шардов (fanout) и доли просматриваемых документов (объём опроса); стоимость хранения --- из дублирования и доли корпуса в крупнейшем шарде. Из правила хранения следует нижняя граница равномерности: шарды терма $t$ содержат все его документы, поэтому
\begin{equation}
	\mathrm{vol}(s) \ge \max_{t:\,s\in\sigma(t)}\mathrm{df}(t),
	\quad
	\mathrm{dup}=\frac{1}{|\mathcal{D}|}\sum_{d\in\mathcal{D}}|D(d)|.
	\label{eq:lower-bound}
\end{equation}
Частые слова ограничивают достижимую равномерность снизу независимо от способа кластеризации. Задача --- определить $\sigma$, доставляющее максимум качеству поиска при заданных ограничениях на перечисленные ресурсы.

\subsection{Граф совстречаемости слов}

Вершины графа --- нормализованные слова обучающих запросов: нижний регистр, стоп-слова, порог документной частоты от~2 до половины корпуса. Слова вопросной рамки (29~позиций: what, mean, difference и~т.\,п.) исключены из графа: они совстречаются с большинством слов и не различают тематику. Ребро соединяет пару с числом совместных появлений не менее двух; вес ребра --- NPMI~\cite{bouma2009}:
\begin{equation}
	\mathrm{NPMI}(a,b)=\ln\frac{n(a,b)\,N}{n(a)\,n(b)}\Bigm/\Bigl(-\ln\frac{n(a,b)}{N}\Bigr),
	\label{eq:npmi}
\end{equation}
где $n(a)$, $n(b)$ --- частоты слов, $n(a,b)$ --- частоты пары, $N$ --- число появлений. Альтернативные веса $G^2$~\cite{dunning1993} и $\chi^2$ и вариации порогов проверены в разделе~5; качественные выводы от них не зависят.

Граф покрывает 5,6\,\% словаря корпуса (10\,567 слов, 39\,794 ребра, 289~компонент связности); размещение остальных слов определяют стратегии шардирования.

\subsection{Построение кластеров}

Кластеризация графа выполнена четырьмя алгоритмами. Leiden~\cite{traag2019} максимизирует модулярность с разрешением $\gamma$:
\begin{equation}
	Q_\gamma=\frac{1}{2m}\sum_{i,j}\Bigl[A_{ij}-\gamma\,\frac{k_i k_j}{2m}\Bigr]\,\delta(c_i,c_j).
	\label{eq:modularity}
\end{equation}
CPM~\cite{reichardt2006} оптимизирует качество constant Potts model с разрешением $\gamma=0{,}02$; Infomap~\cite{rosvall2008} минимизирует ожидаемую длину кода случайного блуждания; METIS~\cite{karypis1998} разрезает граф на 341 часть. Разрешение Leiden $\gamma=1$ и остальные параметры алгоритмов зафиксированы до измерений (раздел~4.4). Алгоритм кластеризации --- заменяемый компонент метода: он задаёт разбиение только слов графа, а распределение нагрузки по шардам определяется стратегиями следующего раздела.

\subsection{Шардирование документов на основе кластеров}

Стратегия \strategy{base} оставляет слова графа в шардах своих кластеров, а вне-графовые слова ($94{,}4\,\%$ словаря) назначает хешем в отдельный диапазон идентификаторов. Стратегия \strategy{aff} присоединяет вне-графовое слово кластеру голосованием документов: берётся кластер со значимым превышением доли совместных документов~\eqref{eq:vote}, где $v(t,c)$ --- число документов со словом $t$ и словами кластера $c$, $m(c)$ --- объём кластера; слова без значимых связей назначаются хешем.
\begin{equation}
	c^*(t)=\arg\max_{c}\;\operatorname{sgn}\Bigl(\frac{v(t,c)}{v(t)}-\frac{m(c)}{m}\Bigr)\,G^2(t,c).
	\label{eq:vote}
\end{equation}
Реплики \strategy{aff\_r2} и \strategy{aff\_r3} дополнительно размещают слово графа в $k$ чужих кластерах наибольшего сродства
\begin{equation}
	\mathrm{aff}(t,c)=\sum_{(t,u)\in E,\;u\in c} w_{tu},
	\label{eq:affinity}
\end{equation}
что сокращает fanout ценой дублирования. Стратегии \strategy{hash\_base}, \strategy{hash\_aff}, \strategy{hash\_bal} --- случайные разбиения на то же число шардов, что у соответствующих семейств; они служат базовой линией.

Стратегия \strategy{bal} строится на разбиении \strategy{aff} и вводит бюджет объёма: доля корпуса в шарде не превосходит $\beta$ (в измерениях $\beta=15\,\%$). Сначала образуются единицы упаковки: кластер в пределах бюджета не дробится; кластер сверх бюджета рекурсивно делится тем же алгоритмом по собственным рёбрам с присоединением вне-графовых слов по значимости совместных документов, пока куски не войдут в бюджет; единица без структуры графа распадается на одиночные слова. Затем единицы, в порядке убывания объёма, упаковываются в шарды с минимальным приростом объёма при соблюдении бюджета:
\begin{equation}
	s^*(u)=\arg\min_{s:\,\mathrm{vol}(s)+\Delta(u,s)\le\beta|\mathcal{D}|}\Delta(u,s),
	\label{eq:packing}
\end{equation}
где $\Delta(u,s)=\bigl|\mathcal{D}(u)\setminus\mathcal{D}(s)\bigr|$. Прирост минимален там, где документы единицы уже хранятся в шарде, поэтому упаковка объединяет единицы с пересекающимися документами и сдерживает рост дублирования; кластер в бюджете делит шард с другими единицами. Слово с $\mathrm{df}(t)>\beta|\mathcal{D}|$ образует шард независимо от остальных. Число шардов --- следствие бюджета, а не параметр.

\subsection{Маршрутизация запроса}

Кандидаты опроса --- шарды термов запроса. Их порядок задаёт жадное взвешенное покрытие~\cite{chvatal1979}: на каждом шаге берётся шард, закрывающий наибольший суммарный вес непокрытых термов,
\begin{equation}
	s_i=\arg\max_{s\in S}\sum_{t\in q\setminus C_{i-1},\; s\in\sigma(t)} w(t),
	\label{eq:routing}
\end{equation}
где $w(t)=\ln\bigl(|\mathcal{D}|/(\mathrm{df}(t)+1)\bigr)+1$ --- IDF терма~\cite{robertson2009}: сначала опрашиваются шарды редких слов, определяющих ранжирование. Бюджет опроса $b$ ограничивает число шардов; $b$, равное мощности покрытия, даёт полный опрос. Алгоритм детерминирован: при равных весах порядок задаётся номером шарда.

\section{Методика оценки}

\subsection{Датасет}

Корпус собран из MS~MARCO~v2.1~\cite{bajaj2016}: релевантные фрагменты обучающих запросов дополнены дистракторами до 241\,185 документов. Запросы (122\,063) разделены в порядке датасета на обучающие (115\,959) и отложенные (6\,104) в отношении 95:5. Словарь корпуса --- 189\,841 слово, 26,2 терма на документ. Граф совстречаемости построен по обучающим запросам и содержит 10\,567 слов и 39\,794 ребра; крупнейшая компонента связности объединяет 9\,972 слова, компонент --- 289.

\subsection{Оценочные выборки}

Основная выборка --- отложенные запросы. Обучающая используется для контроля переобучения графа. Смещённая выборка OOD --- 10\,000 вопросов Natural Questions~\cite{kwiatkowski2019}, распределение которых отличается от построившего граф. Синтетические выборки по 2\,000 запросов зависят только от графа и общи для всех алгоритмов: пары и тройки слов, встречающихся почти только вместе (connected), и наиболее частые пары и тройки (frequent). Число термов запроса --- до пяти.

\subsection{Метрики}

Эталон выдачи --- первые десять документов полного индекса, ранжированного BM25 ($k_1=1{,}2$, $b=0{,}75$)~\cite{robertson2009}. Для каждого бюджета опроса $b\in\{1,2,3,5,\text{полный}\}$ измеряются: доля эталонных документов, найденных в опрошенных шардах (overlap),
\begin{equation}
	\mathrm{overlap}@k(q,b)=\frac{\bigl|R_k(q)\cap\mathcal{D}(S_q^{(b)})\bigr|}{k};
	\label{eq:metrics}
\end{equation}
доля корпуса, просмотренная в опрошенных шардах (объём опроса); fanout --- число шардов, покрывающих термы запроса; доля запросов, покрываемых одним шардом; дублирование; доля корпуса в крупнейшем шарде. Средние приводятся с 95\,\%-ми доверительными интервалами (бутстреп по запросам, 1\,000 перевыборок); сравнения стратегий --- как парные разности на одних и тех же запросах.

Расчёт выдачи шардов выполнен эмуляцией по выдаче полного индекса и сверен с физическими индексами Whoosh, построенными для шардов покрытий проверочных запросов: сопоставлено 1\,594\,686 оценок, максимальное расхождение $3{,}6\cdot10^{-15}$, расхождений выдачи 0. Дальнейшие измерения используют эмуляцию.

\subsection{Сравниваемые конфигурации}

Измерены десять стратегий (\strategy{hash\_base}, \strategy{base}, \strategy{base\_r2}, \strategy{base\_r3}, \strategy{hash\_aff}, \strategy{aff}, \strategy{aff\_r2}, \strategy{aff\_r3}, \strategy{hash\_bal}, \strategy{bal}) на четырёх алгоритмах кластеризации. Основная конфигурация зафиксирована до измерений и не менялась: бюджет $\beta=15\,\%$, разрешение Leiden $\gamma=1$, разрешение CPM $\gamma=0{,}02$, число частей METIS 341, зерно 42, вес рёбер NPMI, вес покрытия IDF. Каждая стратегия сравнивается со случайным разбиением на то же число шардов; чувствительность к зерну, объёму журнала и параметрам алгоритмов проверена отдельно.

\section{Результаты}

\subsection{Сравнение стратегий}

\begin{table*}[!t]
	\ntvcaption{Стратегии шардирования на отложенной выборке, алгоритм Leiden}
	{Sharding strategies on the held-out sample, Leiden}
	\label{tab:strategies}
	\centering
	\resizebox{\textwidth}{!}{%
	\begin{tabular}{|l|r|r|r|r|r|r|r|r|r|r|}
		\hline
		\textbf{Стратегия} & \textbf{Шардов} & \textbf{Дубл.} & \textbf{Макс. шард, \%} & \textbf{Fanout} & \textbf{В 1 шард, \%} & \multicolumn{2}{c|}{\textbf{Overlap@10}} & \multicolumn{3}{c|}{\textbf{Объём опроса, \%}} \\
		 &  &  &  &  &  & @1 & @2 & @1 & @2 & полн. \\
		\hline
		\strategy{hash\_base} & 680 & 25,7 & 15,7 & 3,23 & 6,6 & 0,651 & 0,954 & 3,9 & 8,0 & 15,2 \\
		\hline
		\strategy{base} & 680 & 16,6 & 79,4 & 2,61 & 13,8 & 0,796 & 0,982 & 35,0 & 67,4 & 94,1 \\
		\hline
		\strategy{base\_r2} & 680 & 19,3 & 95,9 & 2,31 & 19,6 & 0,839 & 0,989 & 52,1 & 89,2 & 110,3 \\
		\hline
		\strategy{base\_r3} & 680 & 21,1 & 98,4 & 2,18 & 23,2 & 0,852 & 0,990 & 58,4 & 92,6 & 108,7 \\
		\hline
		\strategy{hash\_aff} & 400 & 25,3 & 19,8 & 3,22 & 6,6 & 0,660 & 0,957 & 6,5 & 12,9 & 23,2 \\
		\hline
		\strategy{aff} & 400 & 10,0 & 83,1 & 2,31 & 22,7 & 0,920 & 0,997 & 59,5 & 104,5 & 136,5 \\
		\hline
		\strategy{aff\_r2} & 400 & 12,2 & 96,4 & 1,98 & 31,3 & 0,952 & 0,999 & 75,7 & 126,8 & 149,4 \\
		\hline
		\strategy{aff\_r3} & 400 & 13,8 & 98,6 & 1,85 & 36,5 & 0,962 & 0,999 & 82,4 & 132,6 & 150,1 \\
		\hline
		\strategy{hash\_bal} & 132 & 23,6 & 31,6 & 3,20 & 6,7 & 0,697 & 0,966 & 17,9 & 35,3 & 60,1 \\
		\hline
		\strategy{bal} & 132 & 19,8 & 15,0 & 2,90 & 9,4 & 0,780 & 0,979 & 15,0 & 28,5 & 43,4 \\
		\hline
	\end{tabular}}
\end{table*}

Семейства без бюджета покупают качество объёмом опроса и размером крупнейшего шарда (табл.~\ref{tab:strategies}). Стратегия \strategy{aff} даёт overlap@10 $0{,}920$ при опросе одного шарда против $0{,}651$ у случайного разбиения, но её крупнейший шард хранит 83,1\,\% корпуса, первый опрос просматривает 59,5\,\% корпуса, а суммарный объём десяти крупнейших шардов --- $6{,}7$ корпуса. Дублирование при этом минимально среди всех стратегий: 10,0 против 25,7 у случайного разбиения --- совстречаемые слова одного кластера, и документы, их содержащие, попадают в меньшее число шардов.

Вырождение устойчиво к настройкам. Среди 18~вариаций порогов графа, весов рёбер, стоп-листа и правила голосования ни одна не приводит крупнейший шард \strategy{aff} к 15-процентному бюджету ни при одном алгоритме; минимум по вариантам --- $20{,}3\,\%$. Ужесточение разрешения Leiden с $\gamma=1$ до $\gamma=20$ сокращает крупнейший шард до $18{,}4\,\%$, но снижает overlap@10 до $0{,}746$ и повышает дублирование до $21{,}0$: параметр кластеризации не заменяет бюджет объёма.

Стратегия \strategy{bal} --- единственная, укладывающаяся в бюджет. При равном числе шардов (132) она превосходит случайное разбиение \strategy{hash\_bal} по всем измеренным метрикам: overlap@10 при опросе одного шарда 0,780 против 0,697 (парная разность +0,083, [0,077;\,0,089]), объём опроса 15,0\,\% против 17,9\,\%, fanout 2,90 против 3,20, дублирование 19,8 против 23,6, крупнейший шард 15,0\,\% против 31,6\,\%. На смещённой выборке OOD прирост составляет +0,109~[0,104;\,0,114]. Выбор устойчив: изменение разрешения Leiden от~1 до~50 меняет overlap@10 \strategy{bal} в пределах $0{,}763$--$0{,}780$ при неизменном объёме шарда; смена зерна кластеризации --- не более чем на $0{,}01$.

\subsection{Выбор конфигурации под профили нагрузки}

Измеренная сетка конфигураций $G$ (алгоритм, стратегия, бюджет опроса) допускает формальный выбор под ограничения эксплуатации. Для каждой конфигурации $c$ известны качество $Q_b(c)$, fanout $F(c)$, дублирование $D(c)$, доля крупнейшего шарда и объём опроса при каждом бюджете. Профиль нагрузки $p$ задаёт метрику $M_p$ и её предельный уровень $M_p^{*}$: минимальный fanout, минимальное дублирование, предельная доля корпуса в шарде, предельный объём опроса. Конфигурация выбирается по правилу
\begin{equation}
	c^{*}(p)=\arg\max_{c\,\in\,G_p} Q_b(c),\qquad
	G_p=\bigl\{c\in G:\; M_p(c)\le M_p^{*},\; Q_b(c)\ge Q_b(H(c))\bigr\},
	\label{eq:selection}
\end{equation}
где $H(c)$ --- случайное разбиение на то же число шардов; условие $Q_b(c)\ge Q_b(H(c))$ отсекает конфигурации, при которых семантическое разбиение не окупается. Равные по качеству конфигурации различаются меньшим fanout и дублированием.

\begin{table*}[!t]
	\ntvcaption{Конфигурации, выбираемые правилом~\eqref{eq:selection} для четырёх профилей нагрузки, и случайное разбиение (отложенная выборка)}
	{Configurations selected by rule~\eqref{eq:selection} for four workload profiles, and the random partition (held-out sample)}
	\label{tab:profiles}
	\centering
	\begin{tabular}{|l|l|r|r|r|r|r|}
		\hline
		\textbf{Профиль} & \textbf{Конфигурация} & \textbf{Fanout} & \textbf{Дублир.} & \textbf{Крупн. шард, \%} & \textbf{Объём опр., \%} & \textbf{Overlap@10} \\
		\hline
		Минимальный fanout & Leiden, \strategy{aff\_r3}, $b=1$ & 1,85 & 13,8 & 98,6 & 82,4 & 0,962 \\
		\hline
		Минимальное дублирование & Leiden, \strategy{aff}, $b=1$ & 2,31 & 10,0 & 83,1 & 59,5 & 0,920 \\
		\hline
		Равномерные шарды & Leiden, \strategy{bal}, $\beta=15\,\%$, $b=1$ & 2,90 & 19,8 & 15,0 & 15,0 & 0,780 \\
		\hline
		Фиксированный объём опроса & CPM, \strategy{aff}, полное покрытие & 2,95 & 21,6 & 32,0 & 14,2 & 1,000 \\
		\hline
		\multicolumn{2}{|l|}{\textit{случайное разбиение (\strategy{hash\_base})}} & 3,23 & 25,7 & 15,7 & 3,9 & 0,651 \\
		\hline
	\end{tabular}
\end{table*}

Применение правила к измеренной сетке даёт конфигурации табл.~\ref{tab:profiles}. Профиль минимального fanout активирует реплики: \strategy{aff\_r3} снижает fanout до~1,85, запросы из совстречающихся слов целиком обслуживаются одним шардом, а качество достигает $0{,}962$; плата --- крупнейший шард $98{,}6\,\%$ и дублирование~13,8. Профиль минимального дублирования выбирает \strategy{aff}: дублирование 10,0 составляет $0{,}39$ от случайного разбиения; плата --- крупнейший шард 83,1\,\% и объём опроса 59,5\,\%. Профиль равномерных шардов реализуется только стратегией \strategy{bal}: доля крупнейшего шарда равна бюджету по построению, и никакая иная конфигурация сетки ограничению не удовлетворяет. Профиль фиксированного объёма опроса при допуске 15\,\% выбирает CPM с \strategy{aff} и полным покрытием: мелкие шарды CPM позволяют просмотреть всё покрытие за $14{,}2\,\%$ корпуса и получить точную выдачу при fanout~2,95; при допуске менее $5\,\%$ правило возвращает мелкошардовые случайные разбиения --- нижняя граница применимости семантического шардирования по этой метрике.

Правило~\eqref{eq:selection} и соответствие «ограничение эксплуатации --- параметр метода» (fanout --- реплики, дублирование --- кластерная структура, равномерность --- бюджет, объём опроса --- размер шардов и бюджет опроса) не зависят от корпуса. Численные значения табл.~\ref{tab:profiles} посчитаны для одного корпуса и одного журнала; перенос на другой корпус требует пересборки сетки тем же конвейером. Гарантии переносимы без пересчёта: доля крупнейшего шарда не превосходит $\beta$ по построению, а сам $\beta$ ограничен снизу частотой самого частого слова --- 11,4\,\% для рассматриваемого корпуса.

\section{Ограничения}

Ранжирование предполагает глобальную статистику BM25 (частоты документов и среднюю длину), доступную всем шардам; локальная статистика изменила бы выдачу и потребовал бы отдельной оценки. Дублирование порядка двадцати копий --- следствие правила полноты~\eqref{eq:invariant}: документ хранится в шардах всех своих термов; сокращение возможно только ценой потери гарантии совпадения выдачи. Совместная встречаемость измерена по журналу запросов: при холодном старте или малом журнале граф беднее --- сокращение журнала со 116 до 10~тыс. запросов снижает overlap@10 \strategy{aff} с $0{,}920$ до $0{,}837$ и повышает дублирование с $10{,}0$ до $14{,}2$. Оценка стоимости выполнена прокси-метриками (число шардов, доля просмотренного корпуса), а не измеренными задержками; нагрузочные эффекты параллельного опроса не моделировались. Все значения получены на одном корпусе и одном языке.

\section{Заключение}

Предложен метод шардирования поискового индекса по словам на основе графа совстречаемости слов журнала запросов с бюджетом объёма шарда. Метод охватывает четыре заменяемых алгоритма кластеризации, голосование документов для вне-графовых слов, реплики по сродству и жадную IDF-взвешенную маршрутизацию; корректность эмуляции подтверждена сверкой с физическими индексами.

Измерения на корпусе MS~MARCO из 241\,185 документов дают три результата. Разбиения без бюджета объёма вырождаются: крупнейший шард достигает 83,1\,\% корпуса, и ни настройка кластеризации, ни 18~вариаций графа не приводят его к бюджету. Кластерная структура сокращает дублирование до 10,0 против 25,7 у случайного разбиения --- главная выгода семантического шардирования. Сбалансированная стратегия \strategy{bal} при равном числе шардов превосходит случайную по всем метрикам: overlap@10 при опросе одного шарда 0,780 против 0,697, на смещённой выборке прирост +0,109.

Сформулировано правило выбора конфигурации под четыре профиля нагрузки; показано, что каждое ограничение эксплуатации отображается на один параметр метода. Дальнейшая работа: перебор бюджета $\beta$ и его влияние на границы профилей, реплики сбалансированной стратегии для совместного профиля равномерности и fanout, измерение задержек на действующей системе и проверка на других корпусах и языках.

\end{ntvbody}

% ============================================================================
\begin{ntvreferences}
\ntvref{kulkarni2015}{Kulkarni A., Callan J. Selective search: Efficient and effective search of large textual collections // ACM Transactions on Information Systems. 2015. V. 33. N 4. P. 17:1–17:33. \url{https://doi.org/10.1145/2738035}}{Kulkarni A., Callan J. Selective search: Efficient and effective search of large textual collections. \textit{ACM Transactions on Information Systems}, 2015, vol.~33, no.~4, pp.~17:1–17:33. \url{https://doi.org/10.1145/2738035}}
\ntvref{dai2016}{Dai Z., Xiong C., Callan J. Query-biased partitioning for selective search // Proceedings of the 25th ACM International Conference on Information and Knowledge Management (CIKM). 2016. P. 1119–1128. \url{https://doi.org/10.1145/2983323.2983706}}{Dai Z., Xiong C., Callan J. Query-biased partitioning for selective search. \textit{Proceedings of the 25th ACM International Conference on Information and Knowledge Management (CIKM)}, 2016, pp.~1119–1128. \url{https://doi.org/10.1145/2983323.2983706}}
\ntvref{moffat2007}{Moffat A., Webber W., Zobel J., Baeza-Yates R. A pipelined architecture for distributed text query evaluation // Information Retrieval. 2007. V. 10. N 3. P. 205–231. \url{https://doi.org/10.1007/s10791-006-9014-4}}{Moffat A., Webber W., Zobel J., Baeza-Yates R. A pipelined architecture for distributed text query evaluation. \textit{Information Retrieval}, 2007, vol.~10, no.~3, pp.~205–231. \url{https://doi.org/10.1007/s10791-006-9014-4}}
\ntvref{risvik2013}{Risvik K.M., Chilimbi T., Tan H. et al. Maguro, a system for indexing and searching over very large text collections // Proceedings of the 6th ACM International Conference on Web Search and Data Mining (WSDM). 2013. P. 727–736. \url{https://doi.org/10.1145/2433396.2433486}}{Risvik K.M., Chilimbi T., Tan H. et al. Maguro, a system for indexing and searching over very large text collections. \textit{Proceedings of the 6th ACM International Conference on Web Search and Data Mining (WSDM)}, 2013, pp.~727–736. \url{https://doi.org/10.1145/2433396.2433486}}
\ntvref{moffat2006}{Moffat A., Webber W., Zobel J. Load balancing for term-distributed parallel retrieval // Proceedings of the 29th Annual International ACM SIGIR Conference on Research and Development in Information Retrieval. 2006. P. 348–355. \url{https://doi.org/10.1145/1148170.1148232}}{Moffat A., Webber W., Zobel J. Load balancing for term-distributed parallel retrieval. \textit{Proceedings of the 29th Annual International ACM SIGIR Conference on Research and Development in Information Retrieval}, 2006, pp.~348–355. \url{https://doi.org/10.1145/1148170.1148232}}
\ntvref{bouma2009}{Bouma G. Normalized (pointwise) mutual information in collocation extraction // Proceedings of the Biennial GSCL Conference. 2009. P. 31–40.}{Bouma G. Normalized (pointwise) mutual information in collocation extraction. \textit{Proceedings of the Biennial GSCL Conference}, 2009, pp.~31–40.}
\ntvref{williams2022}{Williams C.K.I. On suspicious coincidences and pointwise mutual information // Neural Computation. 2022. V. 34. N 10. P. 2037–2046. \url{https://doi.org/10.1162/neco_a_01533}}{Williams C.K.I. On suspicious coincidences and pointwise mutual information. \textit{Neural Computation}, 2022, vol.~34, no.~10, pp.~2037–2046. \url{https://doi.org/10.1162/neco_a_01533}}
\ntvref{dunning1993}{Dunning T. Accurate methods for the statistics of surprise and coincidence // Computational Linguistics. 1993. V. 19. N 1. P. 61–74. \url{https://aclanthology.org/J93-1003}}{Dunning T. Accurate methods for the statistics of surprise and coincidence. \textit{Computational Linguistics}, 1993, vol.~19, no.~1, pp.~61–74. \url{https://aclanthology.org/J93-1003}}
\ntvref{newman2004}{Newman M.E.J., Girvan M. Finding and evaluating community structure in networks // Physical Review E. 2004. V. 69. N 2. P. 026113. \url{https://doi.org/10.1103/PhysRevE.69.026113}}{Newman M.E.J., Girvan M. Finding and evaluating community structure in networks. \textit{Physical Review E}, 2004, vol.~69, no.~2, pp.~026113. \url{https://doi.org/10.1103/PhysRevE.69.026113}}
\ntvref{traag2019}{Traag V.A., Waltman L., van Eck N.J. From {Louvain} to {Leiden}: guaranteeing well-connected communities // Scientific Reports. 2019. V. 9. N 1. P. 5233. \url{https://doi.org/10.1038/s41598-019-41695-z}}{Traag V.A., Waltman L., van Eck N.J. From {Louvain} to {Leiden}: guaranteeing well-connected communities. \textit{Scientific Reports}, 2019, vol.~9, no.~1, pp.~5233. \url{https://doi.org/10.1038/s41598-019-41695-z}}
\ntvref{reichardt2006}{Reichardt J., Bornholdt S. Statistical mechanics of community detection // Physical Review E. 2006. V. 74. N 1. P. 016110. \url{https://doi.org/10.1103/PhysRevE.74.016110}}{Reichardt J., Bornholdt S. Statistical mechanics of community detection. \textit{Physical Review E}, 2006, vol.~74, no.~1, pp.~016110. \url{https://doi.org/10.1103/PhysRevE.74.016110}}
\ntvref{traag2011}{Traag V.A., Van Dooren P., Nesterov Y. Narrow scope for resolution-limit-free community detection // Physical Review E. 2011. V. 84. N 1. P. 016114. \url{https://doi.org/10.1103/PhysRevE.84.016114}}{Traag V.A., Van Dooren P., Nesterov Y. Narrow scope for resolution-limit-free community detection. \textit{Physical Review E}, 2011, vol.~84, no.~1, pp.~016114. \url{https://doi.org/10.1103/PhysRevE.84.016114}}
\ntvref{rosvall2008}{Rosvall M., Bergstrom C.T. Maps of random walks on complex networks reveal community structure // Proceedings of the National Academy of Sciences. 2008. V. 105. N 4. P. 1118–1123. \url{https://doi.org/10.1073/pnas.0706851105}}{Rosvall M., Bergstrom C.T. Maps of random walks on complex networks reveal community structure. \textit{Proceedings of the National Academy of Sciences}, 2008, vol.~105, no.~4, pp.~1118–1123. \url{https://doi.org/10.1073/pnas.0706851105}}
\ntvref{karypis1998}{Karypis G., Kumar V. A fast and high quality multilevel scheme for partitioning irregular graphs // SIAM Journal on Scientific Computing. 1998. V. 20. N 1. P. 359–392. \url{https://doi.org/10.1137/S1064827595287997}}{Karypis G., Kumar V. A fast and high quality multilevel scheme for partitioning irregular graphs. \textit{SIAM Journal on Scientific Computing}, 1998, vol.~20, no.~1, pp.~359–392. \url{https://doi.org/10.1137/S1064827595287997}}
\ntvref{chvatal1979}{Chv{\'a}tal V. A greedy heuristic for the set-covering problem // Mathematics of Operations Research. 1979. V. 4. N 3. P. 233–235. \url{https://doi.org/10.1287/moor.4.3.233}}{Chv{\'a}tal V. A greedy heuristic for the set-covering problem. \textit{Mathematics of Operations Research}, 1979, vol.~4, no.~3, pp.~233–235. \url{https://doi.org/10.1287/moor.4.3.233}}
\ntvref{robertson2009}{Robertson S., Zaragoza H. The probabilistic relevance framework: {BM25} and beyond // Foundations and Trends in Information Retrieval. 2009. V. 3. N 4. P. 333–389. \url{https://doi.org/10.1561/1500000019}}{Robertson S., Zaragoza H. The probabilistic relevance framework: {BM25} and beyond. \textit{Foundations and Trends in Information Retrieval}, 2009, vol.~3, no.~4, pp.~333–389. \url{https://doi.org/10.1561/1500000019}}
\ntvref{bajaj2016}{Bajaj P., Campos D., Craswell N. et al. {MS MARCO}: A human generated {MAchine Reading COmprehension} dataset // arXiv preprint arXiv:1611.09268. 2016. \url{https://doi.org/10.48550/arXiv.1611.09268}}{Bajaj P., Campos D., Craswell N. et al. {MS MARCO}: A human generated {MAchine Reading COmprehension} dataset. \textit{arXiv preprint arXiv:1611.09268}, 2016. \url{https://doi.org/10.48550/arXiv.1611.09268}}
\ntvref{kwiatkowski2019}{Kwiatkowski T., Palomaki J., Redfield O. et al. Natural Questions: a benchmark for question answering research // Transactions of the Association for Computational Linguistics. 2019. V. 7. P. 453–466. \url{https://doi.org/10.1162/tacl_a_00276}}{Kwiatkowski T., Palomaki J., Redfield O. et al. Natural Questions: a benchmark for question answering research. \textit{Transactions of the Association for Computational Linguistics}, 2019, vol.~7, pp.~453–466. \url{https://doi.org/10.1162/tacl_a_00276}}
\end{ntvreferences}

\begin{ntvauthorsinfo}
\ntvauthor{\textbf{Постнов Степан Андреевич}~--- студент, Университет ИТМО, Санкт-Петербург, 197101, Российская Федерация, stepaha78@gmail.com \tbd{ORCID}}
{\textbf{Stepan A. Postnov}~--- Student, ITMO University, Saint Petersburg, 197101, Russian Federation, stepaha78@gmail.com}
\ntvauthor{\textbf{Платонов Алексей Владимирович}~--- кандидат технических наук, доцент, доцент, Университет ИТМО, Санкт-Петербург, 197101, Российская Федерация, https://orcid.org/0000-0002-8485-1296, avplatonov@itmo.ru}
{\textbf{Alexey V. Platonov}~--- PhD, Associate Professor, ITMO University, Saint Petersburg, 197101, Russian Federation, https://orcid.org/0000-0002-8485-1296, avplatonov@itmo.ru}
\ntvauthor{\textbf{Саськов Лев Кириллович}~--- Университет ИТМО, Санкт-Петербург, 197101, Российская Федерация \tbd{должность, ORCID, e-mail}}
{\textbf{Lev K. Saskov}~--- ITMO University, Saint Petersburg, 197101, Russian Federation}
\end{ntvauthorsinfo}

\end{document}
